\قسمت{احراز هویت}

فرض کنید وب‌سایت شما از \متن‌لاتین{Digest Authentication} استفاده می‌کند، مهاجم بسته‌ای شامل اطلاعات احراز هویت کاربر را به سرقت برده است و دوباره آن را برای سرور شما ارسال می‌کند. آیا با اینکار احراز هویت اتفاق می‌افتد؟
اگر پاسخ شما خیر است، در رابطه با روشی که باعث می‌شود چنین چیزی رخ ندهد توضیح دهید.
اگر وب‌سایت شما بر پایه \متن‌لاتین{Basic Authentication} عمل می‌کرد چطور؟

\نمره{۲}

% پاسخ پیشنهادی است و پیش از استفاده در تصحیح نیاز به بازبینی دارد.
\begin{پاسخ}

در \متن‌لاتین{Digest Authentication} سرور در هر بار یک \متن‌لاتین{nonce} تازه می‌دهد و پاسخ کلاینت درهم‌سازی رمز عبور به همراه همان \متن‌لاتین{nonce} است.
بنابراین بسته‌ی ضبط‌شده تنها تا زمانی معتبر است که همان \متن‌لاتین{nonce} پذیرفته شود؛ اگر سرور \متن‌لاتین{nonce} را یک‌بارمصرف کند (یا شمارنده‌ی \متن‌لاتین{nc} را بررسی کند) تکرار بسته احراز هویت نمی‌شود.

در \متن‌لاتین{Basic Authentication} کل نام کاربری و رمز عبور به شکل \متن‌لاتین{base64} در هر تقاضا فرستاده می‌شود.
اینجا تکرار بسته همیشه کار می‌کند و بدتر از آن، مهاجم خودِ رمز عبور را هم دارد.

\شروع{فقرات}
\فقره عدم اشاره به \متن‌لاتین{nonce} به عنوان دلیل کار نکردن تکرار: از دست دادن ۵۰ درصد نمره
\فقره عدم اشاره به لو رفتن رمز عبور در \متن‌لاتین{Basic}: از دست دادن ۲۵ درصد نمره
\پایان{فقرات}

\end{پاسخ}
